\documentclass[12pt,a4paper]{article}
\usepackage[T2A]{fontenc}
\usepackage[utf8]{inputenc}
\usepackage[russian]{babel}
\usepackage{geometry}
\usepackage{amsmath,amssymb,mathtools}
\usepackage{array,booktabs,longtable}
\usepackage{xcolor}
\usepackage{hyperref}
\usepackage{tikz}
\geometry{left=22mm,right=19mm,top=23mm,bottom=23mm}
\hypersetup{colorlinks=true,linkcolor=blue!55!black,urlcolor=blue!55!black}
\setlength{\parindent}{0pt}
\setlength{\parskip}{0.55em}
\renewcommand{\arraystretch}{1.25}
\newcommand{\subpart}[1]{\par\medskip\noindent\textbf{#1}\par\nobreak\smallskip}
\newcommand{\rate}{\text{ тыс. руб.}}
\begin{document}
\begin{titlepage}
\centering
\vspace*{0.28\textheight}
{\Huge\bfseries Ревью домашней работы\par}
\vspace{2.5cm}
{\large Автор: Лёвин Артём Александрович\par}
\vspace{0.7cm}
{\large 20 октября 2026 года\par}
\vfill
\begin{tikzpicture}
\draw[blue!35!black,line width=0.7pt]
  (-4,0) .. controls (-2.3,0.65) and (-1.3,-0.55) .. (0,0)
  .. controls (1.4,0.55) and (2.3,-0.65) .. (4,0);
\end{tikzpicture}
\end{titlepage}

\section*{Сводная таблица}
\small
\begin{longtable}{@{}p{0.85cm}p{2.65cm}p{5.8cm}p{2.15cm}p{2.2cm}@{}}
\toprule
\textbf{№} & \textbf{Тема} & \textbf{Суть недочёта} & \textbf{Ваш ответ} & \textbf{Верный ответ} \\
\midrule
\endhead
\hyperref[task:5]{5} & Два равных платежа & Расчёт остатка долга записан, но не приравнен нулю; первоначальная сумма не найдена. & Не указан & $360\rate$ \\
\addlinespace
\hyperref[task:6]{6} & Три равных платежа & После третьего начисления не вычтен платёж и не найден первоначальный долг. & Не указан & $455\rate$ \\
\addlinespace
\hyperref[task:7]{7} & Размер ежегодного платежа & Решение и ответ отсутствуют на предоставленных фотографиях. & Не представлен & $216\rate$ \\
\addlinespace
\hyperref[task:8]{8} & Неравные платежи & Решение и ответ отсутствуют на предоставленных фотографиях. & Не представлен & $400\rate$ \\
\bottomrule
\end{longtable}
\normalsize

\clearpage
\thispagestyle{empty}
\vspace*{\fill}
\begin{center}
{\LARGE\bfseries Разбор ошибок}
\end{center}
\vspace*{\fill}

\clearpage
\phantomsection\label{task:5}
\section*{Задание №\,5}
\subpart{Текст задания}
Кредит погашается двумя равными ежегодными платежами по $250$ тысяч рублей.
Годовая ставка составляет $25\%$. Найдите первоначальную сумму кредита.

\subpart{Ваше решение}
В обозначениях $S$ --- первоначальная сумма кредита (в тысячах рублей)
Вы составили последовательность остатков долга:
\[
D_0=S,\qquad D_1=1{,}25S-250,
\]
\[
D_2=1{,}25(1{,}25S-250)-250.
\]
Это правильное начало: сначала увеличивается остаток долга на $25\%$,
затем из него вычитается ежегодный платёж.

\subpart{Ваш ответ}
Не указан.

\subpart{Ошибка}
\textcolor{red}{Ошибка:} решение осталось незавершённым.
В приведённых преобразованиях неверного действия нет, однако последнее
выражение не превращено в уравнение, поэтому искомое $S$ не найдено.

\subpart{Где именно возникает недочёт}
Последний правильный шаг --- запись долга после второго платежа:
\[
D_2=1{,}25(1{,}25S-250)-250.
\]
Следующий обязательный шаг отсутствует. Поскольку по условию кредит
выплачен полностью после второго платежа, нужно записать $D_2=0$.
Нельзя считать сумму найденной, пока не решено полученное уравнение.

\subpart{Почему это важно}
Выражение $D_2$ само по себе описывает возможный остаток кредита при
неизвестном $S$. В нём ещё не использовано решающее условие задачи:
\emph{после второго платежа долга не остаётся}. Именно это условие
позволяет найти единственное исходное значение кредита. Платёж $250$
тысяч рублей вычитается \emph{после} начисления процентов, а не до него.

\subpart{Ключевая идея}
\textcolor{blue!60!black}{Ключевая идея:} на каждом шаге сначала
умножаем оставшийся долг на $1{,}25$, потом вычитаем $250$;
после последнего платежа остаток равен нулю.

\subpart{Исправление от последнего правильного шага}
Продолжим именно Вашу запись:
\begin{align*}
1{,}25(1{,}25S-250)-250&=0,\\
1{,}5625S-312{,}5-250&=0,\\
1{,}5625S-562{,}5&=0,\\
1{,}5625S&=562{,}5,\\
S&=\frac{562{,}5}{1{,}5625}=360.
\end{align*}

\subpart{Полное правильное решение}
Пусть первоначальный кредит составляет $S$ тысяч рублей.
Через год долг увеличится на четверть и составит $1{,}25S$.
После первого платежа останется
\[
D_1=1{,}25S-250.
\]
За второй год проценты начисляются не на первоначальную сумму,
а на этот оставшийся долг. Непосредственно перед вторым платежом сумма
к возврату равна $1{,}25D_1$. После платежа останется
\[
D_2=1{,}25D_1-250
   =1{,}25(1{,}25S-250)-250.
\]
Поскольку этот платёж последний, $D_2=0$. Раскрывая скобки и решая
линейное уравнение, получаем $S=360$.

\subpart{Проверка}
При $S=360$ тысяч рублей после начисления первых процентов будет
$360\cdot1{,}25=450$ тысяч рублей. После платежа останется
$450-250=200$ тысяч рублей. Через второй год долг вырастет до
$200\cdot1{,}25=250$ тысяч рублей; второй платёж полностью его погасит:
$250-250=0$. Условия выполнены.

\subpart{Итог}
\textcolor{teal!60!black}{Итог:} первоначальный кредит равен
\boxed{360\rate}.

\subpart{Задача на закрепление}
Кредит погашается двумя равными ежегодными платежами по $180$ тысяч рублей.
Годовая ставка составляет $20\%$. Найдите первоначальную сумму кредита.

\clearpage
\phantomsection\label{task:6}
\section*{Задание №\,6}
\subpart{Текст задания}
Кредит погашается тремя равными ежегодными платежами по $216$ тысяч рублей.
Годовая ставка составляет $20\%$. Найдите первоначальную сумму кредита.

\subpart{Ваше решение}
Вы верно начали составлять остатки долга (суммы в тысячах рублей):
\begin{align*}
D_0&=S,\\
D_1&=1{,}2S-216,\\
D_2&=1{,}2(1{,}2S-216)-216
     =1{,}44S-475{,}2.
\end{align*}
Для третьего года Вы записали начисление процентов на сумму
$1{,}44S-475{,}2$, но окончательный остаток и уравнение не записали.

\subpart{Ваш ответ}
Не указан.

\subpart{Ошибка}
\textcolor{red}{Ошибка:} расчёт прекращён до вычитания третьего платежа
и использования условия полного погашения кредита.

\subpart{Где именно возникает недочёт}
Последний правильно найденный остаток после второго платежа:
\[
D_2=1{,}44S-475{,}2.
\]
Верно начать третий год выражением $1{,}2D_2$, но это долг
\emph{перед} третьим платежом. Следует ещё вычесть $216$ и получить
\[
D_3=1{,}2(1{,}44S-475{,}2)-216=0.
\]

\subpart{Почему это важно}
Проценты начисляются на остаток долга \emph{в начале года}.
Умножение на $1{,}2$ отражает лишь начисление процентов;
само по себе оно не означает, что очередной платёж уже внесён.
Условие задачи предусматривает ровно три платежа, поэтому после
третьего остаток должен равняться нулю. Без этого равенства
известной суммы кредита не получить.

\subpart{Ключевая идея}
\textcolor{blue!60!black}{Ключевая идея:} повторяем правило
«начислить $20\%$ --- вычесть $216$» ровно три раза, а последний
остаток приравниваем нулю.

\subpart{Исправление от последнего правильного шага}
\begin{align*}
D_3&=1{,}2(1{,}44S-475{,}2)-216,\\
0&=1{,}728S-570{,}24-216,\\
1{,}728S&=786{,}24,\\
S&=\frac{786{,}24}{1{,}728}=455.
\end{align*}

\subpart{Полное правильное решение}
Обозначим первоначальный кредит через $S$ тысяч рублей.
В первый год начисляется $20\%$, то есть долг увеличивается
в $1{,}2$ раза. После первого платежа:
\[
D_1=1{,}2S-216.
\]
Во второй год проценты начисляются уже на $D_1$:
\begin{align*}
D_2&=1{,}2D_1-216\\
   &=1{,}2(1{,}2S-216)-216\\
   &=1{,}44S-259{,}2-216\\
   &=1{,}44S-475{,}2.
\end{align*}
В третий год после начисления процентов и последнего платежа
долга не остаётся:
\begin{align*}
0&=1{,}2D_2-216\\
 &=1{,}2(1{,}44S-475{,}2)-216\\
 &=1{,}728S-786{,}24.
\end{align*}
Поэтому $S=455$ тысяч рублей.

\subpart{Проверка}
После первого года: $455\cdot1{,}2-216=330$ тысяч рублей.
После второго года: $330\cdot1{,}2-216=180$ тысяч рублей.
После третьего года: $180\cdot1{,}2-216=0$.
Кредит полностью погашен третьим платежом.

\subpart{Итог}
\textcolor{teal!60!black}{Итог:} первоначальная сумма кредита равна
\boxed{455\rate}.

\subpart{Задача на закрепление}
Кредит погашается тремя равными ежегодными платежами по $125$ тысяч рублей.
Годовая ставка составляет $25\%$. Найдите первоначальную сумму кредита.

\clearpage
\phantomsection\label{task:7}
\section*{Задание №\,7}
\subpart{Текст задания}
Кредит $455$ тысяч рублей взят под $20\%$ годовых и погашается
тремя равными ежегодными платежами. Найдите величину каждого платежа.

\subpart{Ваше решение}
На доступных фотографиях решение задания №\,7 не представлено.

\subpart{Ваш ответ}
Не представлен.

\subpart{Ошибка}
\textcolor{red}{Ошибка:} в присланных материалах отсутствует решение.
Проверить отдельные вычислительные шаги невозможно, поэтому
ни один конкретный неверный переход Вам не приписывается.

\subpart{Где именно возникает недочёт}
Последний корректный шаг Вашего решения установить нельзя:
оно не показано. Недостающий этап --- выразить остаток долга после
трёх равных платежей через неизвестную сумму платежа и приравнять
этот остаток нулю.

\subpart{Почему без этого нельзя получить ответ}
Начальная сумма кредита известна, но неизвестен одинаковый платёж.
На остаток каждый год начисляется $20\%$, а платёж вносится
после начисления процентов. Нельзя просто разделить $455$ на три:
дополнительно выплачиваются проценты, зависящие от того,
какая часть долга ещё не погашена.

\subpart{Ключевая идея}
\textcolor{blue!60!black}{Ключевая идея:} обозначим один ежегодный
платёж через $A$ и трижды применим формулу
«остаток после платежа = остаток до года, умноженный на $1{,}2$,
минус $A$».

\subpart{Исправление}
Поскольку представленного решения нет, начинаем с введения
неизвестного платежа $A$ (в тысячах рублей). После третьего платежа
долг должен стать равным нулю:
\[
1{,}2\bigl(1{,}2(1{,}2\cdot455-A)-A\bigr)-A=0.
\]

\subpart{Полное правильное решение}
После начисления процентов за первый год до платежа долг равен
\[
455\cdot1{,}2=546.
\]
После первого платежа остаётся $D_1=546-A$.
Второй год даёт
\begin{align*}
D_2&=1{,}2D_1-A\\
   &=1{,}2(546-A)-A\\
   &=655{,}2-2{,}2A.
\end{align*}
После третьего начисления процентов и третьего платежа:
\begin{align*}
D_3&=1{,}2D_2-A\\
   &=1{,}2(655{,}2-2{,}2A)-A\\
   &=786{,}24-2{,}64A-A\\
   &=786{,}24-3{,}64A.
\end{align*}
Так как кредит погашен, $D_3=0$. Следовательно,
\begin{align*}
3{,}64A&=786{,}24,\\
A&=\frac{786{,}24}{3{,}64}=216.
\end{align*}

\subpart{Проверка}
Первый год: $455\cdot1{,}2-216=330$ тысяч рублей.
Второй год: $330\cdot1{,}2-216=180$ тысяч рублей.
Третий год: $180\cdot1{,}2-216=0$.
Все три платежа равны и кредит погашен полностью.

\subpart{Итог}
\textcolor{teal!60!black}{Итог:} размер каждого ежегодного платежа
равен \boxed{216\rate}.

\subpart{Задача на закрепление}
Кредит $488$ тысяч рублей взят под $25\%$ годовых и погашается
тремя равными ежегодными платежами. Найдите величину каждого платежа.

\clearpage
\phantomsection\label{task:8}
\section*{Задание №\,8}
\subpart{Текст задания}
Кредит взят на три года под $10\%$ годовых.
Первые два платежа равны по $150$ тысяч рублей, третий ---
$185{,}9$ тысячи рублей. После третьего платежа кредит погашен
полностью. Найдите первоначальную сумму кредита.

\subpart{Ваше решение}
На доступных фотографиях решение задания №\,8 не представлено.

\subpart{Ваш ответ}
Не представлен.

\subpart{Ошибка}
\textcolor{red}{Ошибка:} решение в присланных материалах отсутствует.
Конкретных ошибочных математических действий в Вашей работе
по этому заданию установить нельзя.

\subpart{Где именно возникает недочёт}
Последний корректный шаг Вашего решения неизвестен, поскольку
решение не представлено. Для выполнения задания необходимо поочерёдно
учесть два платежа по $150$ тысяч рублей, затем заключительный
платёж $185{,}9$ тысячи рублей и условие нулевого остатка.

\subpart{Почему требуется учитывать порядок платежей}
Проценты за каждый следующий год начисляются на сумму,
оставшуюся \emph{после предыдущего платежа}. Поэтому нельзя
просто сложить все три платежа и принять их сумму за первоначальный
кредит: часть выплат составляет проценты. Размер третьего платежа
отличается от первых двух, так что в последнем шаге нужно
вычесть именно $185{,}9$, а не $150$.

\subpart{Ключевая идея}
\textcolor{blue!60!black}{Ключевая идея:} трижды увеличиваем
остаток долга на $10\%$, после каждого года вычитаем соответствующий
платёж, а после третьего приравниваем остаток нулю.

\subpart{Исправление}
Поскольку решения на фотографиях нет, начинаем с неизвестного $S$
(тысяч рублей) и записываем полное погашение:
\[
1{,}1\bigl(1{,}1(1{,}1S-150)-150\bigr)-185{,}9=0.
\]

\subpart{Полное правильное решение}
После начисления процентов за первый год и первого платежа остаётся
\[
D_1=1{,}1S-150.
\]
После второго года проценты начисляются на $D_1$:
\begin{align*}
D_2&=1{,}1D_1-150\\
   &=1{,}1(1{,}1S-150)-150\\
   &=1{,}21S-165-150\\
   &=1{,}21S-315.
\end{align*}
В третий год начисляем проценты на $D_2$ и вычитаем последний платёж:
\begin{align*}
D_3&=1{,}1D_2-185{,}9\\
   &=1{,}1(1{,}21S-315)-185{,}9\\
   &=1{,}331S-346{,}5-185{,}9\\
   &=1{,}331S-532{,}4.
\end{align*}
По условию $D_3=0$, значит
\begin{align*}
1{,}331S&=532{,}4,\\
S&=\frac{532{,}4}{1{,}331}=400.
\end{align*}

\subpart{Проверка}
В конце первого года перед платежом долг составляет
$400\cdot1{,}1=440$ тысяч рублей, после платежа остаётся
$440-150=290$ тысяч рублей.
В конце второго года перед платежом долг составляет
$290\cdot1{,}1=319$ тысяч рублей, после платежа остаётся
$319-150=169$ тысяч рублей.
В третьем году долг увеличивается до $169\cdot1{,}1=185{,}9$
тысячи рублей. Последний платёж $185{,}9$ полностью погашает кредит.

\subpart{Итог}
\textcolor{teal!60!black}{Итог:} первоначальная сумма кредита равна
\boxed{400\rate}.

\subpart{Задача на закрепление}
Кредит взят на три года под $20\%$ годовых. Первые два ежегодных
платежа составляют по $200$ тысяч рублей, третий --- $336$ тысяч рублей.
После третьего платежа кредит погашен полностью.
Найдите первоначальную сумму кредита.

\par\vfill
\begin{center}
\footnotesize Лёвин Артём Александрович эксклюзивно для Тимофея Васильченко
\end{center}
\end{document}